\documentclass[12pt]{article}
\usepackage[utf8]{inputenc}
\usepackage[T1]{fontenc}
\usepackage[margin=1in]{geometry}
\usepackage{setspace}
\usepackage{xcolor}
\usepackage[colorlinks=true, urlcolor=blue!60!black, linkcolor=blue!60!black]{hyperref}
\setlength{\parindent}{2em}
\setlength{\parskip}{1.25em}
\renewcommand{\baselinestretch}{1.0}
\title{\textit{Home Building \& Loan Ass'n v. Blaisdell}, 290 U.S. 398 (1934)}
\author{Case Brief}
\date{}
\begin{document}
\maketitle

% Case Brief: Home Building & Loan Ass'n v. Blaisdell, 290 U.S. 398 (1934)
% Constitutional Law Briefing Book — Chapter: Due Process, Contracts & Economic Liberties

\section{Case Metadata}

\begin{description}
  \item[Case Name:] \emph{Home Building \& Loan Ass'n v. Blaisdell}
  \item[Citation:] 290 U.S. 398, 54 S. Ct. 231, 78 L. Ed. 413 (1934) (Docket No.\ 370)
  \item[Dates:] Argued November 8--9, 1933; Decided January 8, 1934
  \item[Court:] Supreme Court of the United States (Appeal from the Supreme Court of Minnesota, \emph{Blaisdell v. Home Building \& Loan Ass'n}, 189 Minn.\ 422, 249 N.W.\ 334 (1933), and 189 Minn.\ 448, 249 N.W.\ 893 (1933))
  \item[Vote:] 5--4, affirming the judgment of the Supreme Court of Minnesota
  \item[Majority Opinion:] Chief Justice Charles Evans Hughes (joined by Justices Louis D.\ Brandeis, Harlan Fiske Stone, Owen J.\ Roberts, and Benjamin N.\ Cardozo)
  \item[Concurring Opinions:] None
  \item[Dissenting Opinion:] Justice George Sutherland (joined by Justices Willis Van Devanter, James C.\ McReynolds, and Pierce Butler)
  \item[Constitutional Provisions:] U.S.\ Const.\ art.\ I, \S\ 10, cl.\ 1 (Contracts Clause); U.S.\ Const.\ amend.\ XIV, \S\ 1 (Due Process and Equal Protection Clauses)
  \item[Primary Opinion Sources:] \href{https://supreme.justia.com/cases/federal/us/290/398/}{Justia U.S.\ Supreme Court Center}; \href{https://www.law.cornell.edu/supremecourt/text/290/398}{Cornell Legal Information Institute (LII)}; \href{https://tile.loc.gov/storage-services/service/ll/usrep/usrep290/usrep290398/usrep290398.pdf}{Library of Congress Official U.S.\ Reports PDF}
\end{description}

\section{Facts and Procedural History}

\subsection{Economic Background and the Minnesota Mortgage Moratorium Law}
At the height of the Great Depression, collapsing farm and urban property values triggered a wave of mortgage defaults and foreclosure sales across the Midwest at severely depressed prices. On April 18, 1933, the Minnesota Legislature enacted the Minnesota Mortgage Moratorium Law, Chapter 339 of the Laws of Minnesota of 1933 (\href{https://www.law.cornell.edu/constitution-conan/article-1/section-10/clause-1/contract-clause}{Cornell LII Constitution Annotated}; \href{https://caselaw.findlaw.com/court/us-supreme-court/290/398.html}{FindLaw}). Reciting the existence of a severe financial and economic depression of several years' duration and the frequent occurrence of foreclosure sales for inadequate prices, the Act declared a public economic emergency calling for the exercise of the State's police power (\href{https://www.law.cornell.edu/constitution-conan/article-1/section-10/clause-1/contract-clause}{Cornell LII}).

Crucially, Chapter 339 did not cancel mortgage debts, strip mortgagees of title rights, or impose a blanket, unconditional freeze. Instead, the statute authorized Minnesota state district courts, upon a mortgagor's application, to extend the statutory period of redemption from mortgage foreclosure sales for such additional time as the court deemed just and equitable, but in no event beyond May 1, 1935 (\href{https://tile.loc.gov/storage-services/service/ll/usrep/usrep290/usrep290398/usrep290398.pdf}{U.S.\ Reports PDF}). As a mandatory condition of any extension, the court was required to determine the reasonable rental value of the property and order the mortgagor remaining in possession to pay all or a reasonable part of that income toward taxes, insurance, interest, and mortgage indebtedness (\href{https://tile.loc.gov/storage-services/service/ll/usrep/usrep290/usrep290398/usrep290398.pdf}{U.S.\ Reports PDF}). The statute also allowed the court to revise or shorten the extension upon changed circumstances, and it declared its relief provisions severable (\href{https://supreme.justia.com/cases/federal/us/290/398/}{Justia}; \href{https://caselaw.findlaw.com/court/us-supreme-court/290/398.html}{FindLaw}).

\subsection{The Underlying Foreclosure and State Court Litigation}
Appellees John H.\ Blaisdell and his wife owned a residential lot in Minneapolis, Minnesota, subject to a mortgage held by appellant Home Building \& Loan Association (\href{https://supreme.justia.com/cases/federal/us/290/398/}{Justia}). After the Blaisdells defaulted, the Association foreclosed on May 2, 1932, under a power of sale contained in the mortgage and purchased the property itself at the foreclosure sale (\href{https://supreme.justia.com/cases/federal/us/290/398/}{Justia}; \href{https://tile.loc.gov/storage-services/service/ll/usrep/usrep290/usrep290398/usrep290398.pdf}{U.S.\ Reports PDF}). Under the Minnesota law in force when the mortgage was executed, the mortgagors had a one-year statutory redemption period, which was set to expire on May 2, 1933.

Following enactment of the Moratorium Law on April 18, 1933, the Blaisdells petitioned the state district court to extend their redemption period. The Association objected to the introduction of evidence and moved to dismiss, contending that Chapter 339 violated the Contracts Clause of Article I, \S\ 10, as well as the Due Process and Equal Protection Clauses of the Fourteenth Amendment (\href{https://caselaw.findlaw.com/court/us-supreme-court/290/398.html}{FindLaw}; \href{https://supreme.justia.com/cases/federal/us/290/398/}{Justia}).

The procedural trajectory unfolded in two rounds before the Minnesota courts:
\begin{enumerate}
  \item \textbf{Initial Dismissal and Reversal:} The trial court initially held the Moratorium Law unconstitutional and dismissed the Blaisdells' petition. The Supreme Court of Minnesota reversed, sustaining the statute as a valid emergency exercise of the State's police power (189 Minn.\ 422, 249 N.W.\ 334 (1933)) (\href{https://supreme.justia.com/cases/federal/us/290/398/}{Justia}).
  \item \textbf{Extension Order and Second Affirmance:} On remand, the trial court found that the reasonable rental value of the premises was \$40 per month and entered a final order extending the redemption period to May 1, 1935, conditioned on the Blaisdells' monthly payment of \$40 to be applied toward taxes, insurance, interest, and mortgage indebtedness (\href{https://tile.loc.gov/storage-services/service/ll/usrep/usrep290/usrep290398/usrep290398.pdf}{U.S.\ Reports PDF}). The Supreme Court of Minnesota affirmed that order (189 Minn.\ 448, 249 N.W.\ 893 (1933)), and Home Building \& Loan Association appealed to the United States Supreme Court (\href{https://supreme.justia.com/cases/federal/us/290/398/}{Justia}).
\end{enumerate}

\section{Issue(s) Presented}

\begin{enumerate}
  \item \textbf{Contracts Clause:} Does a state statute enacted during a severe economic emergency that temporarily authorizes state courts to extend the redemption period for foreclosed mortgages---conditioned upon the mortgagor's payment of court-ascertained reasonable rental value toward carrying costs and interest---violate Article I, \S\ 10, cl.\ 1, which prohibits any State from passing any ``Law impairing the Obligation of Contracts''?
  \item \textbf{Fourteenth Amendment:} Does such a temporary, judicially supervised mortgage moratorium deprive the mortgagee-purchaser of property without due process of law or deny the equal protection of the laws under the Fourteenth Amendment?
\end{enumerate}

\section{Holding}

\textbf{No to both issues.} In a 5--4 decision authored by Chief Justice Hughes, the Supreme Court affirmed the judgment of the Supreme Court of Minnesota and upheld the Minnesota Mortgage Moratorium Law (\href{https://supreme.justia.com/cases/federal/us/290/398/}{Justia}). The Court held that the Contracts Clause is not an absolute prohibition to be applied with literal mathematical exactness; rather, it must be harmonized with the State's reserved police power to safeguard the vital economic interests of its people during an emergency (\href{https://supreme.justia.com/cases/federal/us/290/398/}{Justia}; \href{https://www.law.cornell.edu/constitution-conan/article-1/section-10/clause-1/blaisdell-case-and-state-modifications-to-private-contracts}{Cornell LII}). Because Minnesota's relief was temporary, limited to a genuine economic emergency, addressed to a legitimate public end, and conditioned on reasonable terms protecting the mortgagee's underlying investment through rental payments, it was a constitutionally valid exercise of state protective power (\href{https://supreme.justia.com/cases/federal/us/290/398/}{Justia}). For the same reasons, the Act satisfied Fourteenth Amendment Due Process, and its classification was not arbitrary under the Equal Protection Clause (\href{https://supreme.justia.com/cases/federal/us/290/398/}{Justia}).

\section{Reasoning}

Chief Justice Hughes's majority opinion rests on four interlocking analytical pillars that reshaped twentieth-century Contracts Clause jurisprudence:

\subsection{Emergency and Constitutional Power}
Hughes began by rejecting the notion that an emergency suspends the Constitution: emergency neither creates new constitutional power nor diminishes constitutional restrictions (\href{https://supreme.justia.com/cases/federal/us/290/398/}{Justia}). However, emergency ``may furnish the occasion for the exercise of power'' already possessed by the State (\href{https://supreme.justia.com/cases/federal/us/290/398/}{Justia}). The question, therefore, was not whether Minnesota could override Article I, \S\ 10, but whether the Contracts Clause---properly construed---forbade this particular exercise of the State's preexisting reserved police power.

\subsection{Non-Literal Construction and Implied Police-Power Reservation}
Examining the historical origins of the Contracts Clause---adopted to curb post-Revolutionary state debtor-relief laws that had undermined credit and commerce---Hughes concluded that the Clause is one of the Constitution's broad provisions requiring judicial construction rather than literal, mechanical application (\href{https://supreme.justia.com/cases/federal/us/290/398/}{Justia}; \href{https://www.law.cornell.edu/constitution-conan/article-1/section-10/clause-1/historical-background-on-the-contract-clause}{Cornell LII}). Drawing on earlier police-power decisions such as \emph{Stone v. Mississippi}, 101 U.S.\ 814 (1880), \emph{Manigault v. Springs}, 199 U.S.\ 473 (1905), and the World War I housing cases (\emph{Block v. Hirsh}, 256 U.S.\ 135 (1921); \emph{Marcus Brown Holding Co. v. Feldman}, 256 U.S.\ 170 (1921)), Hughes reasoned that every private contract is made subject to the State's continuing sovereign authority to protect the vital interests of the community (\href{https://supreme.justia.com/cases/federal/us/290/398/}{Justia}; \href{https://www.law.cornell.edu/constitution-conan/article-1/section-10/clause-1/blaisdell-case-and-state-modifications-to-private-contracts}{Cornell LII}).

Rather than relying on rigid formalisms---such as whether a statute affects a contract ``directly'' versus ``incidentally,'' or alters a ``remedy'' versus a ``substantive right'' (acknowledging \emph{Von Hoffman v. City of Quincy}, 71 U.S.\ (4 Wall.) 535 (1867), that validity and remedy are intertwined)---Hughes pivoted to a purposive means--ends inquiry: ``whether the legislation is addressed to a legitimate end and the measures taken are reasonable and appropriate to that end'' (\href{https://supreme.justia.com/cases/federal/us/290/398/}{Justia}). He distinguished nineteenth-century foreclosure cases such as \emph{Bronson v. Kinzie}, 42 U.S.\ (1 How.) 311 (1843), and \emph{Barnitz v. Beverly}, 163 U.S.\ 118 (1896), on the ground that those statutes granted unconditional extensions of redemption without compensating the mortgagee-purchaser for lost possession (\href{https://caselaw.findlaw.com/court/us-supreme-court/290/398.html}{FindLaw}; \href{https://supreme.justia.com/cases/federal/us/290/398/}{Justia}).

\subsection{Evolutionary Constitutional Interpretation}
Anticipating the dissent's originalist critique, Hughes explicitly rejected static originalism. Even though the Framers enacted the Contracts Clause to halt debtor-relief stay laws, Hughes argued that the ``great clauses of the Constitution'' cannot be confined solely to the specific applications the Framers would have envisioned under eighteenth-century conditions; they must be read in light of the Nation's ``whole experience'' and growing economic interdependence (\href{https://supreme.justia.com/cases/federal/us/290/398/}{Justia}; \href{https://www.law.cornell.edu/constitution-conan/article-1/section-10/clause-1/blaisdell-case-and-state-modifications-to-private-contracts}{Cornell LII}).

\subsection{The Five-Factor Balancing Test Applied to Minnesota's Act}
Applying this framework, Hughes identified five specific criteria demonstrating that Minnesota's Moratorium Law stayed within constitutional bounds (\href{https://supreme.justia.com/cases/federal/us/290/398/}{Justia}; \href{https://tile.loc.gov/storage-services/service/ll/usrep/usrep290/usrep290398/usrep290398.pdf}{U.S.\ Reports PDF}):
\begin{enumerate}
  \item \textbf{Existence of an Emergency:} A genuine, severe economic emergency existed in Minnesota, furnishing a proper occasion for invoking the State's protective power.
  \item \textbf{Legitimate Public Purpose:} The legislation was enacted not for the mere advantage of particular individuals, but to protect a basic interest of society---preventing mass homelessness and economic collapse.
  \item \textbf{Appropriately Tailored Relief:} The relief was tailored to the emergency and could be granted only upon reasonable, judicially supervised conditions akin to historic equitable jurisdiction in foreclosure.
  \item \textbf{Protection of the Mortgagee's Core Interest:} The conditions were not unreasonable. The integrity of the mortgage indebtedness remained unimpaired; interest continued to run; the foreclosure sale and right to title or a deficiency judgment upon non-redemption remained intact; and the mortgagor in possession had to pay court-ascertained rental value (\$40/month) toward taxes, insurance, and interest. Hughes also took judicial notice that mortgagees in Minnesota were predominantly corporations concerned with protecting investment security rather than occupying the homes themselves.
  \item \textbf{Temporary Operation:} The legislation was strictly temporary, terminating no later than May 1, 1935, and subject to judicial shortening if circumstances changed. Hughes emphasized that whether the emergency continued to exist remained ``always open to judicial inquiry,'' and the Act could not validly outlast the emergency or virtually destroy contracts.
\end{enumerate}

\subsection{Due Process and Equal Protection}
Hughes disposed of the Fourteenth Amendment claims summarily: the same considerations that established the reasonableness of the moratorium under the Contracts Clause defeated the substantive due process challenge, and limiting relief to mortgage foreclosures was not an arbitrary classification under the Equal Protection Clause (\href{https://supreme.justia.com/cases/federal/us/290/398/}{Justia}).

\section{Concurring and Dissenting Opinions}

\subsection{Concurring Opinions}
No separate concurring opinions were filed.

\subsection{Dissenting Opinion of Justice Sutherland (Joined by Justices Van Devanter, McReynolds, and Butler)}
Writing for the four conservative dissenters (later dubbed the ``Four Horsemen''), Justice Sutherland delivered one of the most celebrated originalist dissents in Supreme Court history (\href{https://supreme.justia.com/cases/federal/us/290/398/}{Justia}; \href{https://www.oxfordreference.com/display/10.1093/oi/authority.20110803095942753}{Oxford Reference}).

\begin{itemize}
  \item \textbf{Fixed Constitutional Meaning:} Sutherland rejected Hughes's evolutionary interpretive method. The meaning of a constitutional provision does not change with economic conditions or shifting public sentiment; its whole aim is to preserve fixed boundaries against passing legislative impulses (\href{https://repository.law.umich.edu/cgi/viewcontent.cgi?article=2353&context=mlr}{Michigan Law Review}; \href{https://digitalcommons.fairfield.edu/cgi/viewcontent.cgi?referer=&httpsredir=1&article=1061&context=nealsb}{Fairfield Digital Commons}).
  \item \textbf{Historical Purpose of the Contracts Clause:} Reviewing the Critical Period of the 1780s (including the wave of state stay laws, installment laws, and tender acts following the Revolutionary War), Sutherland argued that the Contracts Clause was drafted \emph{precisely} to bar state legislatures from enacting debtor-relief moratoria during severe economic depressions. Upholding Minnesota's Act because of an economic depression turned the Clause upside down by suspending its prohibition in the exact circumstance it was designed to govern.
  \item \textbf{Rejection of the Right/Remedy and Compensation Distinctions:} Sutherland contended that under \emph{Bronson v. Kinzie} and \emph{Barnitz v. Beverly}, extending a redemption period after foreclosure directly impairs the contractual obligation by depriving the mortgagee-purchaser of the bargained-for right to immediate title and possession upon expiration of the statutory year. Paying court-set rent was no substitute for the contract's actual terms.
  \item \textbf{Slippery Slope Warning:} Sutherland warned that by replacing a clear constitutional prohibition with an open-ended balancing test dependent on legislative declarations of ``emergency'' and judicial notions of ``reasonableness,'' the majority opened the door to the progressive erosion of constitutional limits whenever they prove politically inconvenient (\href{https://www.scotusblog.com/2026/02/scotustoday-for-thursday-february-12/}{SCOTUSblog}).
\end{itemize}

\section{Key Quotations}

All quotations below are verified verbatim against the official U.S.\ Reports text as published on \href{https://supreme.justia.com/cases/federal/us/290/398/}{Justia} and \href{https://www.law.cornell.edu/supremecourt/text/290/398}{Cornell LII}:

\begin{enumerate}
  \item \textbf{Emergency Does Not Create Power (Chief Justice Hughes, 290 U.S.\ at 425--26):}
  \begin{quote}
    ``Emergency does not create power. Emergency does not increase granted power or remove or diminish the restrictions imposed upon power granted or reserved. The Constitution was adopted in a period of grave emergency. Its grants of power to the federal government and its limitations of the power of the States were determined in the light of emergency, and they are not altered by emergency. \ldots\ While emergency does not create power, emergency may furnish the occasion for the exercise of power.'' (\href{https://supreme.justia.com/cases/federal/us/290/398/}{Justia}; \href{https://www.law.cornell.edu/supremecourt/text/290/398}{Cornell LII})
  \end{quote}

  \item \textbf{Non-Literal Construction of the Contracts Clause (Chief Justice Hughes, 290 U.S.\ at 428):}
  \begin{quote}
    The prohibition against impairment of the obligation of contracts ``is not an absolute one, and is not to be read with literal exactness, like a mathematical formula.'' (\href{https://supreme.justia.com/cases/federal/us/290/398/}{Justia}; quoted in \emph{City of El Paso v. Simmons}, 379 U.S.\ 497, 508 (1965), \href{https://supreme.justia.com/cases/federal/us/379/497/}{Justia})
  \end{quote}

  \item \textbf{Harmonizing the Contracts Clause and Police Power (Chief Justice Hughes, 290 U.S.\ at 439):}
  \begin{quote}
    ``The reserved power cannot be construed so as to destroy the limitation, nor is the limitation to be construed to destroy the reserved power in its essential aspects.'' (\href{https://www.law.cornell.edu/constitution-conan/article-1/section-10/clause-1/blaisdell-case-and-state-modifications-to-private-contracts}{Cornell LII ConAn})
  \end{quote}

  \item \textbf{Evolutionary Interpretation of the Great Clauses (Chief Justice Hughes, 290 U.S.\ at 442--44):}
  \begin{quote}
    ``If by the statement that what the Constitution meant at the time of its adoption it means to-day, it is intended to say that the great clauses of the Constitution must be confined to the interpretation which the framers, with the conditions and outlook of their time, would have placed upon them, the statement carries its own refutation. \ldots\ With a growing recognition of public needs and the relation of individual right to public security, the court has sought to prevent the perversion of the clause through its use as an instrument to throttle the capacity of the States to protect their fundamental interests.'' (\href{https://www.law.cornell.edu/constitution-conan/article-1/section-10/clause-1/blaisdell-case-and-state-modifications-to-private-contracts}{Cornell LII ConAn})
  \end{quote}

  \item \textbf{Constitutional Fidelity When Provisions ``Pinch'' (Justice Sutherland, Dissenting, 290 U.S.\ at 483):}
  \begin{quote}
    ``The only legitimate inquiry we can make is whether it is constitutional. If it is not, its virtues, if it have any, cannot save it; if it is, its faults cannot be invoked to accomplish its destruction. If the provisions of the Constitution be not upheld when they pinch, as well as when they comfort, they may as well be abandoned.'' (\href{https://supreme.justia.com/cases/federal/us/290/398/}{Justia}; \href{https://www.law.cornell.edu/supremecourt/text/290/398}{Cornell LII})
  \end{quote}
\end{enumerate}

\section{Doctrinal Significance and Modern Legacy (Through 2026)}

\subsection{From Immediate Guardrails to Broad Deference}
Initially, the Court enforced \emph{Blaisdell}'s limits, striking down permanent or uncompensated debtor-relief statutes in \href{https://supreme.justia.com/cases/federal/us/292/426/}{\emph{W.\ B.\ Worthen Co.\ v.\ Thomas}, 292 U.S.\ 426 (1934)}, and \href{https://supreme.justia.com/cases/federal/us/295/56/}{\emph{W.\ B.\ Worthen Co.\ v.\ Kavanaugh}, 295 U.S.\ 56 (1935)}. By 1940, however, the Court dropped the emergency requirement in \href{https://supreme.justia.com/cases/federal/us/310/32/}{\emph{Veix v.\ Sixth Ward Building \& Loan Ass'n}, 310 U.S.\ 32 (1940)}, and later upheld a Texas statute cutting off perpetual land-redemption rights in \href{https://supreme.justia.com/cases/federal/us/379/497/}{\emph{City of El Paso v.\ Simmons}, 379 U.S.\ 497 (1965)}.

The Burger Court briefly revived Article I, \S\ 10 in two cases: \href{https://supreme.justia.com/cases/federal/us/431/1/}{\emph{United States Trust Co.\ v.\ New Jersey}, 431 U.S.\ 1 (1977)} (applying heightened ``reasonable and necessary'' scrutiny when a State impairs its \emph{own} financial contract), and \href{https://supreme.justia.com/cases/federal/us/438/234/}{\emph{Allied Structural Steel Co.\ v.\ Spannaus}, 438 U.S.\ 234 (1978)} (invalidating a narrow, retroactive private pension liability statute). In \href{https://supreme.justia.com/cases/federal/us/459/400/}{\emph{Energy Reserves Group, Inc.\ v.\ Kansas Power \& Light Co.}, 459 U.S.\ 400 (1983)}, the Court synthesized a deferential test for private contracts requiring (1) substantial impairment, (2) a significant, legitimate public purpose, and (3) reasonable, appropriate means.

\subsection{\emph{Sveen v.\ Melin} (2018) and Post-Pandemic Circuit Split}
In \href{https://www.law.cornell.edu/supremecourt/text/16-1432}{\emph{Sveen v.\ Melin}, 584 U.S.\ 811 (2018)}, an 8--1 Court streamlined \emph{Blaisdell}/\emph{Energy Reserves} into a two-step test and upheld Minnesota's retroactive revocation-on-divorce statute at step one because a policyholder could undo it with minimal paperwork. Justice Gorsuch dissented alone, echoing Sutherland's originalist critique that the Framers viewed existing contracts as ``inviolable'' (\href{https://www.law.cornell.edu/constitution-conan/article-1/section-10/clause-1/historical-background-on-the-contract-clause}{Cornell LII ConAn}).

Through 2026, the Supreme Court has not struck down a state law under the Contracts Clause since 1978, but COVID-19 moratoria sparked a notable circuit split (\href{https://harvardlawreview.org/print/vol-136/the-contract-clause-reawakened-in-the-age-of-covid-19/}{Harvard Law Review, 2023}): the Ninth Circuit upheld Los Angeles eviction moratoria deferentially (\emph{Apartment Ass'n of L.A.\ County v.\ City of Los Angeles}, 10 F.4th 905 (9th Cir.\ 2021); \href{https://www.govinfo.gov/content/pkg/USCOURTS-ca9-22-55480/pdf/USCOURTS-ca9-22-55480-0.pdf}{\emph{Iten v.\ County of Los Angeles}, 81 F.4th 979 (9th Cir.\ 2023)}), whereas the Second and Eighth Circuits revived challenges to permanent guaranty cancellations (\emph{Melendez v.\ City of New York}, 16 F.4th 992 (2d Cir.\ 2021); \href{https://www.govinfo.gov/content/pkg/USCOURTS-ca2-23-00683/pdf/USCOURTS-ca2-23-00683-0.pdf}{2d Cir.\ 2024}) and indefinite eviction bans (\emph{Heights Apartments, LLC v.\ Walz}, 30 F.4th 720 (8th Cir.\ 2022)) that lacked \emph{Blaisdell}'s temporary, compensated guardrails.

\section{Pedagogical Critiques and Exam Hypotheticals}

\subsection{Critiques and ``Opposite Side'' Contrast}
\emph{Blaisdell} pits living constitutionalism against originalism on a clause expressly aimed at debtor-relief stay laws, raising law-and-economics concerns that ex post relief increases ex ante credit costs and showing how ``emergency'' exceptions normalize into routine deference. Note also that Article I, \S\ 10 binds only States; federal retroactivity is reviewed under Fifth Amendment Due Process (\href{https://supreme.justia.com/cases/federal/us/467/717/}{\emph{PBGC v.\ R.\ A.\ Gray \& Co.}, 467 U.S.\ 717 (1984)}) or the Takings Clause (\href{https://supreme.justia.com/cases/federal/us/295/555/}{\emph{Louisville Joint Stock Land Bank v.\ Radford}, 295 U.S.\ 555 (1935)}).

\textbf{``Opposite Side'' Contrast:} Contrast \emph{Blaisdell} (2-year cap, individualized need, \$40/month rental compensation) with (1) \href{https://supreme.justia.com/cases/federal/us/295/56/}{\emph{Kavanaugh}} (6.5-year foreclosure delay with zero rent or hardship showing---struck down); (2) \href{https://supreme.justia.com/cases/federal/us/438/234/}{\emph{Spannaus}} (permanent, targeted retroactive pension liability---struck down); and (3) \href{https://supreme.justia.com/cases/federal/us/431/1/}{\emph{U.S.\ Trust}} (State impairing its \emph{own} bond covenant triggers heightened ``reasonable and necessary'' scrutiny---struck down).

\subsection{Exam Hypotheticals}
\begin{enumerate}
  \item \textbf{Permanent Commercial Rent Forgiveness:} After a hurricane, State X permanently cancels 100\% of six months' commercial rent and personal guaranties without requiring tenant hardship or compensating landlords. \emph{Analysis:} Fails \emph{Sveen}/\emph{Blaisdell}---substantial impairment destroying core lease terms, and permanent, uncompensated relief mirrors \emph{Kavanaugh} and \emph{Melendez}.
  \item \textbf{Public Pension COLA Repeal:} To close a budget deficit without raising taxes, State Y permanently repeals a statutory 3\% annual COLA for already-retired teachers. \emph{Analysis:} Public contract triggers \emph{U.S.\ Trust} heightened scrutiny; because less drastic fiscal alternatives exist, retroactive repeal is likely neither necessary nor reasonable.
  \item \textbf{Regulated-Industry Notice Rule:} State Z requires solar lessors in a heavily regulated market to mail a one-page notice before automatic 10-year renewals take effect. \emph{Analysis:} Fails at \emph{Sveen} step one---heavy prior regulation lowers expectations (\emph{Energy Reserves}), and minimal paperwork preserves the core bargain.
\end{enumerate}

\end{document}
